\documentclass[11pt]{article}
\usepackage{fontspec}
\usepackage[french]{babel}
\usepackage{amsmath}
\usepackage{amssymb}
\usepackage{amsthm}
\usepackage{mathrsfs}
\usepackage{tikz-cd}
\usepackage{enumitem}

\frenchsetup{AutoSpacePunctuation=false}

\newcommand{\printedpage}[1]{\marginpar{\footnotesize p.~#1}}
\newcommand{\msp}[1]{\marginpar{\footnotesize [#1]}}
\newcommand{\Fibu}{\underline{\mathrm{Fib}}\,}
\newcommand{\Ho}{\mathrm{Ho}}
\newcommand{\WL}{W_S^{\mathcal{L}}}
\newcommand{\defeq}{\overset{\text{déf}}{=}}

\begin{document}

\printedpage{59}
% head: A. Grothendieck : Dérivateurs, Ch. VI, Édition des Universités de Montpellier II et Paris VII
\[
\begin{tikzcd}
\mathcal{F}W_S^{-1} \arrow[r, shift left, "\delta"] & (\mathcal{M}/S)W_S^{-1}, \arrow[l, shift left, "\gamma"]
\end{tikzcd}
\]
il faut voir qu'ils sont quasi-inverse l'une de \msp{page 94} l'autre.

Du côté de $\mathcal{M}/S$ : on doit comparer $X$ et $\delta\gamma(X) = C(X)$, c'est fait par
\[
i_X : X \longrightarrow C(X)
\]
qui est dans $W_S$, d'où l'isomorphisme dans le localisé
\[
X \simeq \delta\gamma(X) \qquad \text{dans } (\mathcal{M}/S)W_S^{-1}.
\]
Il faut vérifier que c'est fonctoriel en $X$.

Du côté de $\mathcal{F}$, comparer $X$ et $\gamma\delta(X)$, i.e. $X$ et $C(X)$. Même tabac, on choisit $i_X : X \simeq C(X) = \gamma\delta(X)$. À vérifier la fonctorialité.

Il s'agit maintenant d'introduire un localisé plus fin de $\mathcal{M}/S$. Pour ceci, on suppose donné une autre classe de morphismes, je l'appelle $\mathcal{L}$ (comme `localisation')
\[
\mathcal{L} \subseteq \mathrm{Fl}(\mathcal{M}),
\]
p. ex. dans le cas de $\mathcal{M} = \mathrm{Cat}$, $\mathcal{L}$ sera formée des isomorphismes locaux (\footnote[55]{Ou aussi, des morphismes \textit{lisses}. Ça doit donner le même résultat.}) (s'il s'agit de définir $W_S^{\mathrm{g}}$), ou des isomorphismes colocaux (\footnote[56]{Ou aussi des morphismes propres.}) (s'il s'agit de définir $W_S^{\mathrm{d}}$). On supposera que $\mathcal{L}$ contient les isomorphismes, et $\mathcal{L}$ stable par composition (\footnote[57]{Il n'en est peut-être pas besoin.}), mais non par changement de base. On définit alors
\[
\WL \subseteq \mathrm{Fl}(\mathcal{M}/S)
\]
par
\[
f \in \WL \overset{\text{déf}}{\Longleftrightarrow} (f \times_S S') \in W,\ \ \forall (S' \longrightarrow S) \in \mathcal{L}.
\]
\msp{page 95}

Comme $\mathcal{L}$ contient $\mathrm{id}_S$, on aura
\[
\WL \subseteq W_S,
\]
d'où un foncteur canonique
\[
\underbrace{(\mathcal{M}/S)(\WL)^{-1}}_{\defeq \Ho_W^{\mathcal{L}}(S)} \xrightarrow{\ \rho\ } \underbrace{(\mathcal{M}/S)(W_S)^{-1}}_{\defeq \Ho_W^{\mathrm{lc}}\mathcal{L}(S)} % [sic] \Ho_W^{\mathrm{lc}}\mathcal{L}(S) -> \Ho_W^{\mathrm{lc}}(S), cf. the definition \Ho_W^{\mathrm{lc}}(S) \defeq (\mathcal{M}/S)W_S^{-1} on p. 61
\]
% foot: Version 9 novembre 2001
\printedpage{60}
% head: A. Grothendieck : Dérivateurs, Ch. VI, Édition des Universités de Montpellier II et Paris VII
qui est un foncteur de localisation. Je voudrais définir un foncteur en sens inverse. Pour ceci, il est naturel de supposer (en analogie avec la situation de Cat)

\textbf{Hypothèse 1.} \textit{Sur $\Fibu S$, $W_S = \WL$, i.e. si $f : X \longrightarrow Y$ est dans $W$ et $X$, $Y$ dans $\Fibu S$, alors pour tout changement de base $S' \longrightarrow S$ qui est dans $\mathcal{L}$, $f' : X_{S'} \longrightarrow Y_{S'}$ est encore dans $W$.}

Dans le contexte de Quillen, cette hypothèse a une nette tendance à être vérifiée quelque soit $\mathcal{L}$, i.e. pout \textit{tout} changement de base $S' \longrightarrow S$ sans restriction. Il suffit en fait que la structure de Quillen soit `propre', plus précisément (et moins restrictivement), qu'un \texttt{[morphisme]} $\in W$ reste tel après changement de base par un morphisme $\in$ Fib. Dégageons ceci comme un Lemme. % [sic] pout -> pour, misprint

\textbf{Lemme.} \textit{L'hypothèse est vérifiée si on suppose} \msp{page 96} \textit{que $F =$ Fib et $W$ sont reliés par la condition supplémentaire (en plus de la factorisation $i_X$) : pour tout carré cartésien}
\[
\begin{tikzcd}
X \arrow[d, "f"'] & X' \arrow[l, "q"'] \arrow[d, "f'"] \\
S & S', \arrow[l, "p"']
\end{tikzcd}
\]
\textit{si $f$ est dans Fib et $p$ dans $W$, alors $q \in W$.}

\textsc{Démonstration}. Soit
\[
\begin{tikzcd}
X \arrow[rr, "f"] \arrow[dr, "p"'] & & Y \arrow[dl, "q"] \\
& S &
\end{tikzcd}
\]
avec $p$, $q \in \mathrm{Fib}$, $f \in W$, montrons que pour \textit{tout} changement de base $g : S' \longrightarrow S$, $f_{S'} : X_{S'} \longrightarrow Y_{S'}$ est encore dans $W$. On factorise $g$ en
\[
\begin{tikzcd}
S' \arrow[r, "i"] \arrow[rr, bend right=15, "g"'] & \bar{S}' \arrow[r, "\bar{g}"] & S,
\end{tikzcd}
\qquad\qquad i \in W,\ g \in \mathrm{Fib},
\]
\texttt{[plutôt} $\bar{g} \in$ \texttt{Fib]} et on est ramené à prouver que l'assertion est vraie dans les deux cas particuliers où $g \in \mathrm{Fib}$, ou $g \in W$. Le cas $g \in \mathrm{Fib}$ résulte aussitôt de l'hypothèse du Lemme (appliquée au carré cartésien
\[
\begin{tikzcd}
X' \arrow[r, "f'"] \arrow[d, "q"'] & Y' \arrow[d, "p"] \\
X \arrow[r, "f"] & Y,
\end{tikzcd}
\]
% foot: Version 9 novembre 2001
\printedpage{61}
% head: A. Grothendieck : Dérivateurs, Ch. VI, Édition des Universités de Montpellier II et Paris VII
où $p \in \mathrm{Fib}$.)

Dans le cas où $S' \longrightarrow S$ est dans $W$, il s'ensuit par l'hypothèse du Lemme appliquée à
\[
\begin{tikzcd}
Y \arrow[d, "\in\mathrm{Fib}"] & Y' \arrow[l, "p"'] \arrow[d] \\
S & S' \arrow[l]
\end{tikzcd}
\]
que \msp{page 97} $p : Y' \longrightarrow Y$ est $\in W$, et de même pour $q : X' \longrightarrow X$, et $f \in W$ implique encore $f' \in W$.

Revenons à l'hypothèse p. 95. On aura donc un foncteur
\[
(\Fibu S, W_S = \WL) \longrightarrow (\mathcal{M}/S, \WL),
\]
d'où pour les localisés un foncteur horizontal
\[
\begin{tikzcd}
(\Fibu S)W_S^{-1} \arrow[r, "i_0"] \arrow[d, "\approx"] & (\mathcal{M}/S)(\WL)^{-1} \defeq \Ho_W^{\mathcal{L}}(S) \arrow[dl] \\
\Ho_W^{\mathrm{lc}}(S) \defeq (\mathcal{M}/S)W_S^{-1} &
\end{tikzcd}
\]
qui factorise le foncteur vertical induit par l'inclusion $(\Fibu S, W_S) \longrightarrow (\mathcal{M}/S, W_S)$ - foncteur dont on a vu (dans les deux cas envisagés) que c'est une équivalence. En inversant ce dernier et le composant avec $i_0$, on trouve donc un foncteur
\[
\underbrace{\Ho_W^{\mathrm{lc}}(S)}_{\defeq (\mathcal{M}/S)W_S^{-1}} \quad \xrightarrow{\ i\ } \quad \underbrace{\Ho_W^{\mathcal{L}}(S)}_{\defeq (\mathcal{M}/S)(\WL)^{-1}} \qquad ,
\]
tel que
\[
\rho i \simeq \mathrm{id},
\]
donc une \textit{rétraction} de $\Ho_W^{\mathcal{L}}(S)$ sur $\Ho_W^{\mathrm{lc}}(S)$. On s'attend que $i$ soit pleinement fidèle, et plus précisément encore, \msp{page 98} que $\rho$ et $i$ soient adjoints, $i$ étant l'adjoint à droite de $\rho$ :
\[
\begin{tikzcd}
\Ho_W^{\mathrm{lc}}(S) \arrow[r, shift right, "i"'] & \Ho_W^{\mathcal{L}}(S), \arrow[l, shift right, "\rho"']
\end{tikzcd}
\]
% foot: Version 9 novembre 2001
\printedpage{62}
% head: A. Grothendieck : Dérivateurs, Ch. VI, Édition des Universités de Montpellier II et Paris VII
i.e. on cherche
\[
\begin{aligned}
i\rho &\quad \underbrace{\xleftarrow{\ i\ }}_{\textit{pas}\ \text{iso}} \quad \mathrm{id} \quad \text{dans } \Ho_W^{\mathcal{L}} \quad (\text{localisation par } \WL)\\
\rho i &\quad \underbrace{\xrightarrow{\ p\ }}_{\text{iso}} \quad \mathrm{id} \quad \text{dans } \Ho_W^{\mathrm{lc}} \quad (\text{localisation par } W_S),
\end{aligned}
\]
\texttt{[trop de} $i$\texttt{'s]} donc pour $X$ dans $\mathcal{M}/S$, on cherche
\[
\begin{aligned}
i_X :&\quad X \quad \longrightarrow \quad C(X) \quad \text{tout au moins dans } (\mathcal{M}/S)(\WL)^{-1}\\
p_X :&\quad X \quad \longleftarrow \quad C(X) \quad \text{tout au moins dans } (\mathcal{M}/S)(W_S)^{-1}.
\end{aligned}
\]
On choisira pour $i_X$ la flèche `canonique' de $X$ dans $C(X)$, définie dans $\mathcal{M}/S$ lui-même, sans avoir à localiser. Mais dans le `cas Quillen' (quand $C$ n'est pas donné comme foncteur $\mathcal{M}/S \longrightarrow \Fibu S$), il faut voir que cet $i_X$ est fonctoriel en $X$, quand on le regarde dans le localisé $(\mathcal{M}/S)W_S^{-1}$. Mais c'est contenu dans la construction de $C$ comme \textit{foncteur}
\[
\mathcal{M}/S \longrightarrow (\Fibu S)W_S^{-1}.
\]
On notera que $i_X$ est dans $W_S$ sans plus, pas dans $\WL$, donc ne définit pas en général un isomorphisme de $\Ho_W^{\mathcal{L}}$.

On définit $p_X$ à l'aide de $i_X$, comme l'inverse de $i_X$ \textit{dans} $(\mathcal{M}/S)W_S^{-1}$, qui existe car \msp{page 99} $i_X$ est dans $W_S$, donc \texttt{[est un]} isomorphisme dans $(\mathcal{M}/S)W_S^{-1} = \Ho_W^{\mathrm{lc}}$, et $p$ est donc \texttt{[un]} isomorphisme. On retrouve donc $\rho i \simeq \mathrm{id}$, donc que $\rho$ est une \textit{rétraction} de $\Ho_W^{\mathcal{L}}$ sur $\Ho_W^{\mathrm{lc}}$.

Pour prouver l'adjonction, en procédant comme dans le cas de Cat, on est ramené à prouver que pour tout $X$,
\[
\left\{
\begin{aligned}
&C(i_X) = i_{C(X)} : C(X) \rightrightarrows C(C(X))\\
&\textit{dans } \Ho_W^{\mathcal{L}} \ (\text{localisation par } \WL).
\end{aligned}
\right.
\]
J'y étais arrivé dans Cat, en utilisant un morphisme
\[
C(C(X)) \xrightarrow{\ c\ } C(X)
\]
(essentiellement, un morphisme de composition des chemins) tel que $c \in \WL$, et que $c \circ C(i_X) = c \circ i_{C(X)}$. Ça avait l'air très délicat, car c'était à grand peine que j'étais arrivé à construire un foncteur \underline{Ch} (`chemins') raisonnable, où on puisse bel et bien composer les chemins! Je renonce pour le moment à axiomatiser les propriétés du foncteur \underline{Ch} dont j'ai besoin pour faire marcher la machine (je devrais y revenir tantôt). Je vais commencer plutôt par le `cas Quillen', avec des $i_X \in \mathrm{Cof}_0$.
% foot: Version 9 novembre 2001

\printedpage{63}
% head: A. Grothendieck : Dérivateurs, Ch. VI, Édition des Universités de Montpellier II et Paris VII
\msp{page 100}

On reprend la description du `foncteur' $C$ (à valeurs dans $(\Fibu S)W_S^{-1}$, et alors sans guillemets), on trouve que $C(i_X)$, où
\[
X \xrightarrow{u = i_X} X' = C(X)
\]
est défini simplement (dans $(\Fibu S)W_S^{-1}$) par `la' flèche $C(u)$ rendant commutatif le diagramme (dans $\Fibu S$)
\[
\begin{tikzcd}
X \arrow[r, "u\ (=i_X)"] \arrow[d, "i_X"'] & X'\ (= C(X)) \arrow[d, "i_{X'}\ (=i_{C(X)})"] \\
C(X) \arrow[r, "C(u)\ (=C(i_X))"] & C(X')\ (= C(C(X))),
\end{tikzcd}
\]
de sorte qu'on a dans $\mathcal{M}/S$ le diagramme
\[
\begin{tikzcd}
X \arrow[r, "\overbrace{i_X}^{\in\mathrm{Cof}_0\cap W}"] & \overbrace{C(X)}^{\in\mathrm{Ob}\,\Fibu S} \arrow[r, shift left, "i_{C(X)}"] \arrow[r, shift right, "C(i_X)"'] & \overbrace{C(C(X))}^{\in\mathrm{Ob}\,\Fibu S},
\end{tikzcd}
\]
commutatif, i.e. $i_{C(X)}i_X = C(i_X)i_X$. Le corollaire page 92 nous dit que les flèches $i_{C(X)}$ et $C(i_X)$ sont égales dans $(\Fibu S)W_S^{-1}$. Mais par hypothèse (p. 95), dans $\Fibu S$ on a $W_S = \WL$, donc les flèches sont égales dans $(\Fibu S)(\WL)^{-1}$, et a fortiori dans $(\mathcal{M}/S)(\WL)^{-1}$, et on gagne.

Voyons ce qu'on peut faire dans le cas où on dispose d'un \textit{foncteur} $C$. On suppose bien sûr que ce foncteur dans $\mathcal{M}/S$ \msp{page 101} n'est pas lié au $S$ particulier dans $\mathcal{M}$ qu'on a choisi comme base, que sa construction provient à la \textsc{Cartan-Serre} d'un foncteur
\[
X \longmapsto \mathrm{Ch}(X) : \mathcal{M} \longrightarrow \mathcal{M},
\]
(`objet des chemins') et de morphismes fonctoriels
\[
X \xrightarrow{\alpha_X} \mathrm{Ch}(X) \xrightarrow{p_X} X \times X, \qquad\qquad \boxed{p_X\alpha_X = \mathrm{diag}_X}
\]
(\footnote[58]{\textbf{NB} On pose $\mathrm{pr}_1 \circ p_X = \sigma_X$, $\mathrm{pr}_2 \circ p_X = \beta_X$ (flèches `source' et `but' pour l'objet $X$).}). Je présume qu'il faudra supposer surtout ceci

\textbf{Hypothèse 2.} \textit{$\alpha_X \in W$, $p_X \in$} Fib.

Pour un $X$ sur $S$, on construira $C_S(X)$ ou $C(X)$ ainsi
% foot: Version 9 novembre 2001

\end{document}
